\documentclass[12pt, a4paper]{article}

% ---------------------------------------------------------------
% PACKAGES — pdflatex + mathpazo (LuaLaTeX fontspec fallback)
% ---------------------------------------------------------------
\usepackage[T1]{fontenc}
\usepackage[utf8]{inputenc}
\usepackage{mathpazo}
\usepackage{microtype}
\usepackage{geometry}
\geometry{
  a4paper,
  top=2.5cm,
  bottom=2.5cm,
  left=3cm,
  right=3cm
}
\usepackage{setspace}
\onehalfspacing
\usepackage{amsmath}
\usepackage{amsthm}
\usepackage{amssymb}
\usepackage{mathtools}
\usepackage{csquotes}
\usepackage[hidelinks]{hyperref}
\usepackage{url}
\usepackage{titlesec}
\titleformat{\section}
  {\large\bfseries}{\thesection.}{1em}{}
\titleformat{\subsection}
  {\normalsize\bfseries}{\thesubsection.}{1em}{}
\usepackage{epigraph}
\setlength{\epigraphwidth}{0.6\textwidth}
\usepackage{booktabs}
\usepackage[hang, flushmargin]{footmisc}

% ---------------------------------------------------------------
% THEOREM ENVIRONMENTS
% ---------------------------------------------------------------
\theoremstyle{plain}
\newtheorem{theorem}{Theorem}[section]
\newtheorem{proposition}[theorem]{Proposition}
\newtheorem{lemma}[theorem]{Lemma}
\newtheorem{corollary}[theorem]{Corollary}

\theoremstyle{definition}
\newtheorem{definition}[theorem]{Definition}
\newtheorem{example}[theorem]{Example}

\theoremstyle{remark}
\newtheorem{remark}[theorem]{Remark}

% ---------------------------------------------------------------
% TITLE
% ---------------------------------------------------------------
\title{
  \textbf{The Geometry of Control}\\[0.5em]
  \large Admissibility, Reachability, and Projection Collapse\\
  in \textit{Prisoners of Power}
}
\author{
  Flyxion\\[0.3em]
  \small Independent Researcher
}
\date{June 2026}

% ---------------------------------------------------------------
% DOCUMENT
% ---------------------------------------------------------------
\begin{document}

\maketitle
\thispagestyle{empty}

\newpage

% ---------------------------------------------------------------
% ABSTRACT
% ---------------------------------------------------------------
\begin{abstract}
Arkady and Boris Strugatsky's \textit{Prisoners of Power}
is commonly interpreted as a political novel concerned with
totalitarianism, propaganda, and the abuse of centralized
authority. While such readings capture important dimensions
of the text, they leave unexplained the novel's persistent
emphasis on uncertainty, hidden structure, failed interventions,
and the repeated collapse of seemingly adequate explanations.
This essay argues that the novel is more fruitfully understood
as an investigation into the geometry of knowledge and control.

Viewed through the lenses of admissibility, reachability,
and projection failure, the narrative becomes a sequence of
increasingly severe epistemic crises. Maxim Kammerer repeatedly
encounters situations in which locally coherent models prove
globally inadequate. Every explanatory framework he constructs
eventually collapses under the weight of distinctions it fails
to preserve. The resulting narrative is not merely political
but geometric: it concerns the relationship between
representations and the futures they permit agents to reach.

At the same time, the novel resists complete assimilation
into any formal framework. Its treatment of loyalty, sacrifice,
historical contingency, and intervention under radical
uncertainty exposes limitations in current theories of
projection and repair. The novel therefore functions not as
an illustration of a theory but as a stress-test for one.
The goal of this essay is not to show that \textit{Prisoners
of Power} confirms a preexisting framework, but to investigate
what becomes visible when the novel and the framework are
allowed to challenge one another.
\end{abstract}

\newpage

% ---------------------------------------------------------------
% TABLE OF CONTENTS
% ---------------------------------------------------------------
\tableofcontents

\newpage

% ---------------------------------------------------------------
% PART I
% ---------------------------------------------------------------
\part*{Part I: The Novel as an Epistemic Machine}
\addcontentsline{toc}{part}{Part I: The Novel as an Epistemic Machine}

% ---------------------------------------------------------------
\section{Introduction: Literature as Theory Discovery}
% ---------------------------------------------------------------

The relationship between literature and theory is often
conceived in one direction. A theoretical framework is
developed independently and subsequently applied to a
literary work. The text becomes an example, a case study,
or a convenient source of illustrations. Such readings
may be useful, but they frequently reduce literature to
a passive object of analysis.

There exists another possibility. Certain literary works
operate not as examples of theories but as instruments
for discovering phenomena that existing theories struggle
to describe. In these cases, the text functions less like
a specimen and more like an experimental apparatus.
Narrative becomes a means of exploring structures that
remain difficult to formalize directly.

Science fiction has historically occupied this role.
Long before formal work on cybernetics, information
theory, artificial intelligence, or complex systems,
science fiction authors were already constructing
environments in which such phenomena could be observed
in exaggerated form. The value of these works does not
arise from predictive accuracy. Rather, it arises from
their ability to isolate structural relationships that
ordinary experience tends to obscure.

The works of Arkady and Boris Strugatsky belong to this
tradition. Although frequently discussed as political
allegories, their novels consistently display a deeper
concern with the relationship between knowledge and action.
Their protagonists rarely confront simple moral dilemmas.
Instead, they encounter worlds whose structure exceeds
the explanatory resources available to them. The resulting
conflicts emerge not primarily from malice or ideology
but from incomplete understanding.

\textit{Prisoners of Power} represents perhaps the clearest
example of this tendency. The novel follows Maxim Kammerer,
a young explorer from the highly developed civilization
of the Noon Universe, after his accidental arrival on
the planet Saraksh. At first glance, the narrative appears
to present a familiar confrontation between freedom and
totalitarian control. Yet this interpretation quickly
proves inadequate. Every major revelation in the novel
destabilizes the explanatory framework that preceded it.
The regime is not what it initially appears to be. The
resistance is not what it initially appears to be. Even
liberation itself becomes difficult to evaluate.

This recurring pattern suggests that the central concern
of the novel is not political authority but epistemic
authority. The fundamental question is not who should
govern. The fundamental question is how agents construct
models of worlds whose most important constraints remain
hidden from view.

Critical interpretation has not always recognized this.
The history of the novel's reception displays a tendency
that the novel itself repeatedly dramatizes at the level
of narrative. Faced with a text that resists familiar
categories, critics have reached for the nearest available
framework. \textit{Prisoners of Power} becomes an example
of dystopian fiction, Soviet allegory, Cold War anxiety,
or anti-totalitarian critique. It is compared to Orwell,
to Lem, to Le Guin, to the tradition of outsider
intervention narratives. These comparisons are not
incorrect. Yet they share a common structure: they
stabilize the text's uncertainty by assigning it to
a known category.

This is precisely what Maxim does on Saraksh.

Faced with an environment whose structure exceeds his
available frameworks, he reaches for the nearest category.
The planet becomes an adventure, then a dystopia, then
a propaganda system, then a liberation project. Each
category captures something real. Each conceals something
essential. The critical reception of the novel reproduces,
at the level of interpretation, the epistemic condition
the novel is investigating at the level of plot.

The present essay attempts a different approach.

Rather than locating the novel within an existing
conceptual landscape, it treats the novel as a source
of theoretical pressure. The question is not which
existing framework best accommodates the text. The
question is what the text reveals that existing frameworks
have not yet reached.

The concepts brought to bear in what follows---projection,
admissibility, reachability, repair---are not applied
to the novel from outside. They emerge as responses to
problems the novel has already made visible. The theory
arrives after the phenomenon, not before it.

This ordering is methodologically significant.

When theory arrives before phenomenon, the phenomenon
becomes evidence for the theory. Confirming instances
are foregrounded. Resistant instances are explained away.
The text is interpreted until it fits.

When theory arrives after phenomenon, the relationship
inverts. The phenomenon becomes a constraint on the
theory. The text is allowed to resist. The places where
the formal vocabulary fails to capture what the narrative
sees are not embarrassments to be minimized. They are
the most valuable results the analysis can produce.

The argument proceeds in six parts. The first three
examine the novel on its own terms, attending to the
specific texture of its epistemic structure before
any formal apparatus is introduced. The middle sections
develop the formal vocabulary and bring it into contact
with the novel's central conflicts. The final sections
examine the limits of the framework and the theoretical
questions the novel generates that the framework cannot
yet answer.

The goal, throughout, is not to reduce the novel to
mathematics. The goal is to place two modes of inquiry
into conversation, and to take seriously the possibility
that the conversation will change both.

% ---------------------------------------------------------------
\section{The Experience of Epistemic Collapse}
% ---------------------------------------------------------------

Most political novels derive their force from revelation.
The protagonist discovers a hidden truth concealed beneath
an official narrative. Once this truth becomes visible,
the moral and political structure of the story becomes
increasingly clear. The reader moves from ignorance
toward understanding.

\textit{Prisoners of Power} operates differently.

Rather than moving from confusion to clarity, the novel
repeatedly moves from local clarity to global uncertainty.
Every revelation generates a larger mystery. Every
explanation succeeds only temporarily before being
displaced by another. The reader does not experience
the accumulation of knowledge so much as the progressive
destabilization of certainty.

This structure is visible from the opening chapters.
Maxim arrives on Saraksh carrying assumptions inherited
from the civilization of the Noon Universe. He expects
contact, cooperation, and rational dialogue. His
understanding of social development encourages him to
interpret unfamiliar institutions as incomplete versions
of familiar ones. The result is not simple
misunderstanding but systematic misprojection. He
continuously interprets local phenomena through
categories whose applicability has not yet been
established.

The reader initially shares these assumptions. As a
consequence, each collapse of Maxim's model
simultaneously becomes a collapse of the reader's model.
The narrative repeatedly reproduces the experience of
discovering that an apparently sufficient description
has omitted distinctions that later prove indispensable.

In this sense, the novel functions as a machine for
generating epistemic humility.

The structure is worth examining in its details, because
the pattern of collapse is not random. Each framework
Maxim adopts is more sophisticated than the one that
preceded it. Each captures more of the available
evidence. Each lasts longer before failing. Yet each
failure is more consequential than the last, because
each collapse occurs after Maxim has already acted on
the basis of the framework that is about to give way.

This asymmetry between the increasing sophistication
of Maxim's models and the increasing severity of their
failures is the novel's primary structural feature.
It is not a coincidence of plot. It reflects a deep
property of the kind of system Saraksh represents:
a system in which increasing local accuracy of
representation does not guarantee increasing global
adequacy, because the most consequential structures
of the system are precisely those that every available
framework is least equipped to preserve.

The reader's experience of this structure is itself
theoretically significant.

One does not merely observe Maxim's failures from
outside. One inhabits them. The novel is constructed
so that the reader's explanatory resources are always
approximately one step behind the situation. The
explanation that seems adequate always turns out to
have omitted something. The framework that seemed to
account for the evidence always turns out to have
collapsed a distinction that the next chapter will
reveal to be essential.

This is not a flaw in the novel's construction.
It is the novel's primary theoretical contribution
at the level of form.

The reader is given the experience, not merely the
description, of operating within a projection that
proves inadequate. The theoretical content is
delivered phenomenologically before it is delivered
propositionally. By the time the formal analysis of
the following chapters arrives, the reader has already
inhabited the condition that analysis is attempting
to describe.

% ---------------------------------------------------------------
\section{The Architecture of Concealment}
% ---------------------------------------------------------------

Before the formal apparatus of projection and
admissibility is introduced, it is necessary to
attend carefully to the specific architecture of
concealment that the novel constructs.

The world of Saraksh is organized as a hierarchy
of hidden constraints, each layer generating the
phenomena that the layer above it misinterprets
as fundamental.

The surface layer is the one Maxim encounters first:
a totalitarian state, perpetually at war, with
omnipresent police, compulsory ideology, and a
population living under conditions of privation and
surveillance. This layer is real. The suffering is
genuine. The coercion is genuine. The institutional
apparatus of authoritarian governance is present
and operative.

But it is not the deepest layer.

Beneath the political surface lies the broadcast
system. The towers scattered across the landscape
are not, as their official designation suggests,
anti-ballistic missile defenses. They are the
primary mechanism through which the cognitive
space of the population is managed. The political
institutions that occupy the surface are in part
downstream effects of this deeper mechanism.
Understanding the surface without understanding
the broadcasts is like understanding a shadow
without understanding the object that casts it.

Beneath the broadcast system lies the distinction
between the immune and the susceptible. The Unknown
Fathers, the degens, and eventually Maxim himself
inhabit a different cognitive space than the
susceptible majority. This distinction organizes
the entire political landscape in ways invisible
to anyone who does not know it exists. The
resistance movement, the structure of governance,
the categories of insider and outsider: all are
shaped by a variable that most inhabitants of the
system cannot observe.

Beneath the immune-susceptible distinction lies
the geopolitical situation that the planet's own
information environment cannot represent. The
Island Empire, the broader context of Saraksh's
history, the presence of Progressors from the
Noon Universe: these are structures operating at
a scale that no native actor's framework is
designed to accommodate.

And beneath all of these lies the physical fact
with which this analysis must begin: the
atmospheric conditions that produce the spherical
world illusion, the most fundamental projection
failure in the novel, the one that precedes every
other distortion and provides the structural
template for all that follows.

% ---------------------------------------------------------------
\section{The First Projection: Living Inside the Sphere}
% ---------------------------------------------------------------

Before the reader encounters the Towers, the Unknown
Fathers, the mutants, or the hidden mechanisms of
political control, the novel introduces a more
fundamental distortion. The inhabitants of Saraksh
believe that they live on the interior surface
of a sphere.

This belief is not merely a cultural superstition.
It emerges from the physical conditions of the
planet itself. Atmospheric refraction prevents
direct observation of the horizon and produces
the appearance of a closed world. The illusion
is sufficiently compelling that it becomes embedded
within the civilization's ordinary understanding
of reality. What is remarkable is not that the
inhabitants are mistaken. What is remarkable is
that the mistake is rational.

The people of Saraksh are not irrational observers
ignoring obvious evidence. On the contrary, they
are responding appropriately to the information
available to them. Their conclusions emerge from
observation, inference, and experience. The error
arises not from defective reasoning but from the
structure of the observations themselves.

This distinction is crucial.

Many accounts of knowledge assume that error
originates primarily from failures of logic.
An individual reaches an incorrect conclusion
because they reason badly. The world itself
remains transparent. The responsibility for
failure lies entirely with the observer.

The situation on Saraksh is different. The
inhabitants inhabit a world whose structure
actively distorts observation. Their reasoning
operates upon projections already shaped by
physical constraints. Under such circumstances,
intelligence alone cannot guarantee correctness.
Better reasoning applied to the same distorted
evidence may simply produce more sophisticated
versions of the same mistake.

The atmospheric illusion therefore functions as
the novel's first and most important lesson.
Before political control appears, before propaganda
appears, before ideological conflict appears, the
reader is confronted with a deeper possibility:
reality may be systematically misrepresented even
in the absence of deception.

The significance of this point extends far beyond
the geography of Saraksh. The spherical-world
illusion provides a miniature version of the
epistemic condition that governs the entire
narrative. Again and again, characters operate
within environments where observations are genuine
yet incomplete, where conclusions are locally
justified yet globally incorrect, and where hidden
structures shape visible phenomena without becoming
directly visible themselves.

The planet itself therefore becomes a metaphor
for representation.

The inhabitants do not see the world directly.
They see a projection generated by interactions
between the world and the medium through which
it is observed. The resulting model successfully
supports many practical activities. People
navigate, build cities, conduct wars, and organize
societies while possessing a fundamentally incorrect
understanding of the larger structure in which
those activities occur.

This possibility should not be dismissed as
merely fictional. The history of science contains
numerous examples of comparable situations.
Ancient astronomy successfully predicted celestial
motions while operating within an incorrect
cosmological framework. Classical mechanics
generated extraordinarily accurate predictions
despite lacking knowledge of relativistic effects.
Early theories of heredity enabled practical
breeding long before the molecular basis of
genetics became understood. In each case, useful
local competence coexisted with incomplete global
understanding.

The lesson is not that knowledge is impossible.
The lesson is that successful action does not
guarantee adequate representation.

The distinction between competence and understanding
becomes one of the novel's recurring themes.
Systems may function while resting upon mistaken
assumptions. Societies may remain stable while
operating under distorted descriptions of reality.
Individuals may act effectively while
misunderstanding the mechanisms that make their
effectiveness possible.

The atmospheric illusion of Saraksh represents
the first instance of this pattern.

It also establishes a structural irony that
persists throughout the novel. The reader encounters
a civilization whose entire worldview is organized
around an unnoticed projection. Yet the reader
simultaneously assumes that this revelation places
them in a privileged epistemic position. Having
recognized the inhabitants' mistake, one is tempted
to believe that the deeper structure of the world
has now become visible.

The remainder of the novel systematically dismantles
this confidence.

Each subsequent revelation reproduces the same
pattern at a different scale. The official ideology
is exposed as a projection. The resistance's
interpretation is exposed as a projection. Maxim's
revolutionary ambitions become a projection. Even
the apparently straightforward opposition between
freedom and control eventually proves inadequate.

The sphere therefore serves not merely as background
worldbuilding but as a structural template for the
entire narrative. Every major conflict in the novel
can be understood as a dispute occurring within a
representational framework that later turns out
to be incomplete.

What changes from episode to episode is not the
existence of projection but the location of the
hidden constraint.

From the perspective of admissibility, the
significance of the spherical illusion becomes
even clearer. The inhabitants possess a
representation that preserves many distinctions
while destroying others. Their model supports
local navigation but fails to preserve the global
geometry of their environment. The representation
is therefore neither wholly wrong nor wholly
correct. It occupies an intermediate position in
which certain forms of action remain possible
while others become systematically distorted.

This observation suggests a more general principle.
Representations should not be evaluated solely in
terms of truth or falsity. They should also be
evaluated according to the futures they preserve
or destroy.

A representation may contain inaccuracies while
remaining highly useful. Conversely, a
representation may contain many correct facts
while obscuring distinctions essential for
navigation. The practical significance of a model
depends not only upon correspondence with reality
but also upon its capacity to preserve reachability.

The civilization of Saraksh survives despite its
mistaken cosmology because the distinctions
required for everyday activity remain largely
intact. The illusion becomes dangerous only when
actions require access to structures that the
projection has erased.

This relationship between projection and
reachability forms the conceptual foundation
of the novel. Long before the Towers begin
broadcasting, long before Maxim becomes involved
in revolutionary politics, the reader has already
been introduced to a world in which representation
and reality have diverged.

The remainder of the narrative can therefore be
read as an exploration of a single question:
what happens when an entire civilization attempts
to navigate using maps that preserve only part
of the territory?

% ---------------------------------------------------------------
\section{Maxim Kammerer and the Confidence of Good Maps}
% ---------------------------------------------------------------

The tragedy of Maxim Kammerer does not arise
from ignorance.

Indeed, Maxim is among the least ignorant
individuals on Saraksh.

He arrives from a civilization possessing
scientific, technological, and social capabilities
far beyond those available on the planet. He
possesses knowledge inaccessible to nearly everyone
he encounters. He is educated, intellectually
curious, courageous, and remarkably resistant to
ideological manipulation. By conventional standards,
he is exactly the kind of individual one would
expect to understand the situation most clearly.

Yet the novel repeatedly demonstrates that Maxim's
advantages become sources of vulnerability.

The reason is subtle. Maxim does not merely possess
knowledge. He possesses confidence in the
relationship between knowledge and understanding.

He assumes that increasing information leads
naturally toward better models. He assumes that
hidden structures, once discovered, can be
incorporated into a progressively more accurate
description of reality. Most importantly, he
assumes that truth and successful intervention
tend to align.

These assumptions are reasonable.

They are also precisely the assumptions the novel
places under pressure.

Throughout the narrative, Maxim repeatedly
discovers facts that overturn previous beliefs.
Each discovery improves his understanding of
some aspect of the situation. Yet the accumulation
of these discoveries does not produce convergence
toward certainty. Instead it reveals additional
layers of complexity. Every solved mystery exposes
a larger unsolved one.

The resulting pattern differs significantly from
traditional heroic narratives. The hero's journey
is often structured around increasing competence.
Obstacles are overcome through learning, adaptation,
and perseverance. Knowledge accumulates and
eventually enables decisive action.

Maxim's trajectory follows the opposite direction.

Knowledge accumulates.
Confidence deteriorates.
The world becomes less legible rather than more.

This distinction is essential for understanding
the novel's deeper structure. Maxim is not
primarily a revolutionary, a soldier, or an
explorer. He is an interpreter. His central
activity throughout the narrative consists of
constructing explanatory models and revising them
when they fail.

The true drama of the novel therefore unfolds
not on battlefields or within political institutions
but within the space of representation itself.

The question is never merely what Saraksh is.

The question is how one determines what Saraksh
is when every available description appears
incomplete.

\subsection{Successive Maps of Saraksh}

The structure of \textit{Prisoners of Power} can
be understood as a succession of maps. Each map
explains more than the one that preceded it. Each
map is also eventually revealed to be inadequate.

The significance of this pattern lies in the fact
that the maps are not simply false. Every
interpretive framework Maxim adopts captures
genuine features of the world. The problem is
never complete error. The problem is partial
adequacy.

This distinction is important because it shifts
attention away from truth and falsity toward
preservation and omission. The central question
becomes not whether a model is correct but whether
it preserves the distinctions upon which future
understanding depends.

Maxim's first map emerges immediately after his
arrival. He interprets Saraksh through the lens
of exploration. The planet appears as an
opportunity for discovery. Its inhabitants seem
likely to become partners in communication and
cultural exchange. The situation resembles the
adventure narratives familiar from earlier eras
of exploration. Maxim imagines himself as a
modern Robinson Crusoe temporarily isolated from
civilization but ultimately destined to establish
productive contact with a new world.

This interpretation is not irrational. The
inhabitants are human. Communication proves
possible. Friendships do emerge. Yet the map
omits the institutional and historical structures
that organize the society Maxim has entered.
The resulting failure is not a consequence of
mistaken observations. It is a consequence of
premature framing. Before enough distinctions
have been gathered, the entire situation has
already been assigned to a category.

The first map collapses.

A second map replaces it. Having encountered
forced labor camps, police surveillance,
ideological rituals, and omnipresent
militarization, Maxim begins to interpret Saraksh
as a totalitarian state. The explanatory structure
now appears considerably stronger. Previously
inexplicable events acquire coherence. The
suffering of ordinary citizens becomes
intelligible. Political authority emerges as the
central organizing principle of social life.

Once again, the interpretation succeeds. The
regime is genuinely authoritarian. Violence is
real. Oppression is real. Yet the map remains
incomplete. The society's visible institutions
are not the primary mechanism through which power
operates. The categories of conventional political
analysis explain certain features of the system
while obscuring others. Maxim has identified an
important layer of reality but mistakes it for
the deepest layer.

The second map collapses.

A third map emerges. The discovery of the
broadcasts appears to provide the missing
explanation. Suddenly the irrational enthusiasm
of the population, the apparent stability of the
regime, and the peculiar status of the degens
become comprehensible. The hidden mechanism has
been revealed.

This stage occupies a privileged position within
the narrative because it resembles the kind of
revelation toward which many political novels
build. The hidden machinery becomes visible.
The protagonist discovers the truth concealed
behind appearances. The reader expects the
trajectory toward liberation to begin.

The novel refuses this structure.

Instead, the discovery of the broadcasts generates
a new simplification. The world becomes divided
into manipulators and manipulated, oppressors and
victims, controllers and controlled. The
explanatory power of this model is considerable.
Yet it also encourages a new reduction: the entire
complexity of the civilization threatens to become
subordinated to a single mechanism.

Here the novel introduces one of its most important
insights. The identification of a hidden mechanism
does not guarantee the identification of an
admissible intervention. Knowledge of cause does
not imply knowledge of consequence. The discovery
of a constraint does not reveal the geometry of
its removal.

The distinction becomes visible through Maxim's
encounters with actors who refuse to share his
certainty. The mutants, in particular, repeatedly
evaluate situations according to criteria different
from his own. Their concern is not whether the
system is unjust. They already know that it is.
Their concern is what futures become reachable
after the proposed repair.

From Maxim's perspective, these responses often
appear passive or overly cautious. From another
perspective, however, they reveal a form of
reasoning largely absent from his own approach.
Whereas Maxim evaluates the present state of the
system, the mutants evaluate the trajectories
emerging from interventions upon it.

The difference is subtle but profound.

One asks what is wrong.
The other asks what follows.

As the narrative progresses, Maxim gradually
acquires additional information regarding the
broader geopolitical situation. The Island Empire,
initially imagined as a potential source of
assistance, proves capable of atrocities comparable
to those committed by the regime he opposes.
The opposition between freedom and tyranny begins
to fracture. Every candidate solution reveals its
own anomalies. Every apparent escape route
terminates in another constraint.

The result is a progressive erosion of explanatory
certainty. What initially appeared to be a contest
between truth and falsehood becomes a contest among
incomplete descriptions. What appeared to be a
struggle between freedom and oppression becomes
a conflict among incompatible futures.

At each stage Maxim's understanding increases.
At each stage the world becomes more difficult
to navigate.

This recurring phenomenon suggests a more precise
characterization of Maxim's failure.

His error is not overconfidence in particular
conclusions. His error is confidence in the
adequacy of the representational frame from which
those conclusions are drawn.

Every map is treated as though it contains the
distinctions necessary to understand the next
decision. Every map eventually proves to have
discarded distinctions that later become
indispensable.

The resulting pattern may be termed projection
confidence: not confidence that one's beliefs
are true, but confidence that the categories
through which beliefs are organized are sufficient.

The tragedy of Maxim Kammerer is therefore not
that he repeatedly reaches incorrect conclusions.
The tragedy is that he repeatedly discovers the
incompleteness of the spaces within which those
conclusions were formed.

The novel's deepest lesson may be that intelligence
alone does not guarantee escape from this condition.
One may continually revise beliefs while remaining
trapped within representational structures whose
limitations become visible only after action has
already begun.

In this sense, Maxim's journey is not a movement
from ignorance to knowledge.

It is a movement from confidence in maps to
awareness of mapping itself.

% ---------------------------------------------------------------
% PART II
% ---------------------------------------------------------------
\newpage
\part*{Part II: Projection Failure}
\addcontentsline{toc}{part}{Part II: Projection Failure}

% ---------------------------------------------------------------
\section{Admissibility, Projection, and Reachability}
% ---------------------------------------------------------------

The preceding chapters have developed a problem
through narrative means.

A civilization rests upon a distorted representation
of its physical environment. A political system
maintains stability through systematic deformation
of cognitive space. An agent with superior knowledge
repeatedly constructs models that prove locally
adequate and globally incomplete. A liberation event
produces consequences indistinguishable, in their
immediate effects, from catastrophe.

These observations share a common structure.

In each case, a representation that appears
sufficient proves to omit distinctions upon which
future navigation depends. In each case, the failure
is not located in reasoning but in the
representational framework within which reasoning
occurs. And in each case, the question of whether
a given representation is adequate cannot be
answered by examining the representation alone.
It requires examining what the representation
preserves and what it destroys.

This section introduces a formal vocabulary for
these phenomena. The concepts are not imported
from outside the problem. They are responses to
the problem that the novel has already forced
into view.

\subsection{State Descriptions and Projection Operators}

Let $\Omega$ denote the space of possible states
of a system. For complex social systems, $\Omega$
is typically high-dimensional, encompassing
institutional configurations, population belief
states, material conditions, historical trajectories,
and latent causal structures invisible to any
particular observer.

No agent has direct access to $\Omega$.

Every agent constructs a representation by selecting
dimensions deemed relevant and discarding others.
This operation can be formalized as a projection
operator $\pi: \Omega \to \Sigma$, where $\Sigma$
is a lower-dimensional representation space.

The projection $\pi$ is not a distortion introduced
by error. It is a necessary feature of finite
cognition. No agent can maintain a complete
description of $\Omega$. Every act of understanding
involves selecting which distinctions to preserve.

The critical question is therefore not whether
projection occurs. It always occurs. The question
is which distinctions survive the projection and
which are discarded.

\subsection{Admissibility}

\begin{definition}
A projection $\pi: \Omega \to \Sigma$ is
\emph{admissible} with respect to a decision class
$\mathcal{D}$ if it preserves the distinctions
necessary for decisions in $\mathcal{D}$ to be
executed without systematic distortion.
\end{definition}

More precisely, let $d \in \mathcal{D}$ be a
decision whose correct execution depends on
distinguishing between states $\omega_1, \omega_2
\in \Omega$. The projection $\pi$ is admissible
for $d$ if $\pi(\omega_1) \neq \pi(\omega_2)$.

If $\pi(\omega_1) = \pi(\omega_2)$ while $\omega_1$
and $\omega_2$ require different responses under
$d$, then $\pi$ is inadmissible for that decision.
The agent will behave identically in situations
that require different behavior. Failure is
geometric rather than logical: the map has
collapsed a distinction the territory requires.

Admissibility is therefore always relative to a
decision class. A projection may be highly
admissible for one class of decisions while being
deeply inadmissible for another.

The civilization of Saraksh provides a clear
illustration. The spherical-world model is
admissible for a wide range of everyday navigational
and social decisions. It becomes inadmissible
precisely when decisions require knowledge of the
planet's true geometry.

\subsection{Reachability}

\begin{definition}
Let $R(\omega, t)$ denote the set of states
accessible from $\omega$ within time horizon $t$
under the available repertoire of actions. This
set constitutes the \emph{reachability volume}
of an agent at state $\omega$.
\end{definition}

A representation $\pi$ preserves reachability
if the actions available within $\Sigma$ can be
mapped back to actions in $\Omega$ without
systematic reduction of $R(\omega, t)$.

A representation fails to preserve reachability
when it collapses distinctions between states
that require different actions to reach the same
target. Under such collapse, the agent believes
a target is reachable when it is not, or believes
it is unreachable when it is, because the
distinctions necessary to identify the correct
action sequence have been discarded.

Reachability failure of this kind is pervasive
in \textit{Prisoners of Power}. Maxim repeatedly
believes that certain futures are reachable from
his current position and acts accordingly. The
mutants' refusals, the Wanderer's interventions,
and the catastrophic aftermath of the Control
Center's destruction all represent situations in
which the reachability volume Maxim believed he
inhabited was larger than the reachability volume
that actually existed.

\subsection{Projection Collapse}

\begin{definition}
\emph{Projection collapse} designates the failure
mode in which a representation that successfully
supports one class of decisions is extended to
a second class for which it was not designed,
with the result that distinctions critical to the
second class are systematically discarded.
\end{definition}

This is the precise structure of Maxim's recurring
error. His exploration frame was constructed for
initial contact decisions. Extended to political
analysis, it collapses distinctions between
institutional types. His political analysis frame
was constructed for understanding surface power
relations. Extended to intervention design, it
collapses distinctions between causal mechanisms.
His broadcast-control frame was constructed for
understanding the mechanism of population
management. Extended to civilizational
intervention, it collapses distinctions between
constraint-removal and stabilization.

In each case, the agent is confident in the
adequacy of a frame constructed for a narrower
purpose than the one to which it is now being
applied. Projection confidence is therefore a
systematic phenomenon rather than an individual
psychological tendency. It arises whenever
successful local application of a representation
generates confidence that the same representation
is adequate for decisions of wider scope.

\subsection{The Broadcast System as Admissibility Architecture}

The Towers of Saraksh can now be characterized
more precisely than the political analysis of
earlier chapters permitted.

The broadcast system is not a content delivery
mechanism. It is an admissibility architecture.

It operates by constraining the projection
operators available to the population. Under
continuous low-intensity broadcast, the
representations through which most citizens
interpret their social environment cannot preserve
certain distinctions. Specifically, they cannot
preserve distinctions between the propaganda model
of social reality and alternative models that
would be constructed by an agent with access to
the suppressed cognitive space.

The twice-daily intense broadcast performs a
complementary function. It does not install
beliefs. It resets affective states in ways that
prevent the accumulation of the dissonance that
would otherwise motivate revision of the available
representational frameworks.

The system therefore operates on the meta-level
of cognition rather than the object level. It
does not determine what citizens conclude. It
determines the space of frameworks within which
conclusions can be formed.

\subsection{Epistemic Admissibility and Civilizational Admissibility}

The analysis now permits a precise formulation
of the novel's central tension.

Maxim optimizes for epistemic admissibility.
He seeks representations that preserve distinctions
between true and false descriptions of reality.
His consistent objective is transparency.

The Wanderer optimizes for what might be termed
civilizational admissibility. He seeks to preserve
the reachability volume of Saraksh as a
civilization across long time horizons. His
consistent objective is navigability.

These objectives are not the same.

A representation may be epistemically inadmissible
while preserving civilizational reachability. The
broadcasts produce systematically distorted models
of reality among the susceptible population. Yet
the civilization continues to function.

A representation may be epistemically admissible
while destroying civilizational reachability.
Sudden access to accurate information about the
actual mechanisms of social control, in the absence
of replacement stabilizing structures, may produce
disorientation, institutional collapse, and
long-term reduction in the reachable futures
available to the civilization.

The Wanderer's argument is precisely that Maxim
has optimized for the wrong admissibility criterion.

Maxim's response is equally coherent. A civilization
maintained through systematic epistemic
inadmissibility is navigating toward futures that
its members have not chosen and cannot evaluate.
The reachability that the Wanderer seeks to preserve
is not the reachability of the population. It is
the reachability of a trajectory selected by a
small immune minority operating under objectives
the majority cannot access.

Both arguments are formally valid.

The novel provides no theorem that resolves them.

What it provides instead is the structure of the
problem: two admissibility criteria, incompletely
correlated, each locally compelling and globally
in tension. The following chapters develop the
formal tools required to state this tension
precisely enough that the question of resolution
can at last be asked clearly.

% ---------------------------------------------------------------
% PART III
% ---------------------------------------------------------------
\newpage
\part*{Part III: The Geometry of Intervention}
\addcontentsline{toc}{part}{Part III: The Geometry of Intervention}

% ---------------------------------------------------------------
\section{Intervention Under Incomplete Observability}
% ---------------------------------------------------------------

Every actor in \textit{Prisoners of Power} who
attempts to change the situation on Saraksh fails.

This observation is worth dwelling upon before
any theoretical apparatus is introduced, because
it is not the kind of observation that conventional
political analysis tends to produce. Political
analysis typically asks which actor had the correct
diagnosis and what prevented them from implementing
it. The implicit assumption is that correct
diagnosis plus sufficient power produces successful
intervention.

The novel systematically refuses this assumption.

The actors who fail include not only those with
incorrect diagnoses but also those with
sophisticated, partially correct, and even
remarkably accurate understandings of the system.
Failure is not distributed according to epistemic
merit. It is distributed according to the structure
of the problem itself.

This suggests that the novel's deepest subject is
not any particular actor's inadequacy. It is the
conditions under which intervention in complex
systems becomes possible at all.

\subsection{Load-Bearing Constraints}

The central error in Maxim's intervention theory
can be stated precisely.

He treats the broadcast system as a pure restriction
on reachable cognitive states. Under this model,
the system reduces the population's accessible
future manifold. Removing the system therefore
expands that manifold. The intervention is
unambiguously beneficial.

The Wanderer's opposing argument rests on a
different structural observation. The broadcast
system is not only a restriction. It is also a
stabilizing constraint on a civilization that has
recently undergone catastrophic nuclear war, that
operates under conditions of severe material
scarcity, that faces external military threats,
and that has organized its social and psychological
life around the rhythms and affective states the
broadcasts produce.

Under these conditions, the constraint does not
merely reduce the population's accessible futures.
It also partially constitutes the navigational
infrastructure within which any future becomes
accessible at all.

The distinction between a restrictive constraint
and a load-bearing constraint is not visible from
within the framework Maxim employs. His framework
preserves the distinction between freedom and
oppression while collapsing the distinction between
constraints that limit and constraints that also
support.

\subsection{The Successive Objective Functions}

The pattern of Maxim's failure across the novel
can be made precise by examining not the content
of his beliefs but the structure of the objective
functions he successively adopts.

In the first phase, the objective is survival and
contact. The frame collapses the distinction
between a society that is less developed in a
neutral sense and one whose institutional structure
is organized around coercion and concealment.

In the second phase, the objective is resistance
to tyranny. The frame captures surface political
organization while collapsing the distinction
between a tyranny maintained through force and
one maintained through neurological manipulation
of the population's representational space.

In the third phase, the objective is elimination
of the broadcast mechanism. The frame correctly
identifies the primary mechanism while collapsing
the distinction between a constraint that limits
and a constraint that also stabilizes.

In the fourth phase, which is largely implicit
in the novel's ending, the objective becomes
post-liberation stabilization. This is the
objective for which Maxim has no prepared
framework at all. He has spent the novel
developing increasingly sophisticated analyses
of what is wrong with Saraksh. He has not
developed any framework for what a navigable
Saraksh would look like or how to reach it from
the current state.

This asymmetry is not accidental. It is a
structural feature of how the novel positions
its protagonist. Maxim's analytical resources
are organized around diagnosis. The problem of
repair under uncertainty receives almost no
theoretical attention until the consequences
of his intervention have already materialized.

\subsection{Why Every Actor Fails}

The systematic failure of every intervention in
the novel suggests that the problem is not located
in any particular actor's limitations. It is
located in the structure of the situation itself.

The Unknown Fathers maintain the broadcast system
as a permanent feature of governance. Their
intervention succeeds in maintaining stability
across decades while systematically preventing
the population from developing the cognitive
resources that would make the control mechanism
unnecessary. The Fathers' intervention preserves
civilizational continuity at the cost of
indefinitely deferring the conditions under which
the intervention could be legitimately ended.

The underground attempts to build resistance
through information and organization. Their
intervention is epistemically admissible in the
sense that it operates through accurate information
and genuine consent. It fails to achieve
civilizational scale because the representational
space of the susceptible majority cannot
accommodate the distinctions the underground's
analysis requires.

Maxim attempts to destroy the control mechanism
directly. His intervention is physically
successful. It produces consequences that exceed
his reachability analysis by an order of magnitude.

The mutants refuse intervention entirely. Their
refusal is epistemically sophisticated: they
correctly identify that the proposed futures
following from available interventions are worse
than the current situation from their perspective.
Yet refusal is itself an intervention in the sense
that it determines which other interventions
become possible.

The Wanderer pursues the most sophisticated
intervention: gradual systemic reform operating
below the threshold of political visibility,
designed to improve conditions incrementally
without triggering the instabilities that rapid
change would produce. His intervention is
destroyed by Maxim's action.

The observation that emerges from surveying these
failures is not that some actors were right and
others wrong. It is that the system resists
intervention at every level of sophistication.

\subsection{The Intervention Problem Formalized}

Let $\mathcal{S}$ denote the state space of a
social system and $\mathcal{I}$ a class of
interventions available to an agent. An
intervention $\iota \in \mathcal{I}$ maps a
current state $s \in \mathcal{S}$ to a
distribution over future states
$\iota(s) \in \Delta(\mathcal{S})$.

The agent's observational access to $\mathcal{S}$
is mediated by a projection $\pi: \mathcal{S}
\to \Sigma$, where $\Sigma$ is the agent's
representational space. The agent selects
interventions based not on $s$ directly but on
$\pi(s)$. If $\pi$ is inadmissible for the
decision class associated with intervention
design, the agent's selection procedure will
systematically produce interventions whose actual
consequences in $\mathcal{S}$ differ from their
anticipated consequences in $\Sigma$.

This gap between anticipated and actual
consequences is not reducible to prediction error
in any ordinary sense. It is a structural
consequence of the mismatch between the
distinctions preserved in $\Sigma$ and the
distinctions upon which the consequences of
interventions in $\mathcal{S}$ depend.

The intervention problem under incomplete
observability can therefore be stated as follows.

Given an agent with projection $\pi$, a system
state $s$, and a target future $s^* \in
\mathcal{S}$, does there exist an intervention
$\iota \in \mathcal{I}$ such that $\iota(s)$
places high probability mass on states near $s^*$,
and can the agent identify this intervention from
within $\Sigma$?

In general, the answer depends on whether the
distinctions required to identify the correct
intervention are preserved by $\pi$. If they are
not, no amount of reasoning within $\Sigma$ will
reliably produce the correct selection.

% ---------------------------------------------------------------
\section{Political Conflict as Competing Repair Programs}
% ---------------------------------------------------------------

The preceding analysis has established that every
actor in \textit{Prisoners of Power} fails to
produce an admissible intervention. What it has
not yet explained is why the actors fail in such
different ways and toward such different ends.

The actors do not simply disagree about methods.
They disagree about what is broken.

\subsection{Anomalies and Repair Operators}

\begin{definition}
Let a system state $s \in \mathcal{S}$ be
evaluated against a reference standard $s^*
\in \mathcal{S}$. An \emph{anomaly} is any
feature of $s$ that diverges from $s^*$ in a
way the evaluating agent regards as requiring
correction. The \emph{anomaly set} $\mathcal{A}$
of an agent is the collection of features they
identify as requiring repair. A \emph{repair
operator} $\rho: \mathcal{S} \to \mathcal{S}$
is an intervention designed to eliminate or
reduce the anomalies in $\mathcal{A}$.
\end{definition}

The anomaly set is not determined by the state
of the system alone. It is determined by the
intersection of the system's current state and
the agent's projection. Features that are
invisible within the agent's projection cannot
appear in their anomaly set. Features that are
visible but assigned low priority may be
structurally important while receiving no
attention from the repair operator. And features
that the agent regards as anomalies may be
load-bearing constraints that the system requires
for stability.

Political conflict, on this analysis, is rarely
a disagreement about facts. It is a collision
between repair operators designed to eliminate
incompatible anomaly sets.

\subsection{The Anomaly Sets of Saraksh}

For the Unknown Fathers, the primary anomaly is
civilizational vulnerability. The broadcasts are
not a problem to be repaired. They are a repair
already implemented.

The underground resistance selects a different
anomaly set. For them, the primary anomaly is
the suppression of cognitive freedom and political
self-determination. The broadcasts are not a
solution but the central problem.

The mutants identify yet another anomaly set.
Their primary concern is survival and the
preservation of their community's distinctive
mode of life. They evaluate proposed repairs
not by asking whether they address the
underground's anomalies or the Fathers' anomalies
but by asking whether the futures following from
those repairs contain viable trajectories for
their own community.

The Wanderer's anomaly set is the most
comprehensive and temporally extended. He
identifies the primary anomaly as the absence
of conditions under which Saraksh can develop
into a self-determining civilization. Every
feature of the current situation is evaluated
against this standard. The broadcasts are
anomalous not because they suppress freedom
in the present but because they prevent the
development of the cognitive and institutional
resources that long-term civilizational health
requires. Yet their sudden removal is equally
anomalous because it destroys the stability
that any developmental trajectory requires
as a foundation.

\subsection{Why Identical Observations Generate
Incompatible Repair Strategies}

When two agents operate from incompatible
projections, the information produced by one
agent's analysis is not straightforwardly
receivable by the other. The distinctions that
make a piece of evidence significant within one
framework may not be preserved within the other.
Argument that proceeds from premises visible
within one projection and invisible within another
will appear, from within the second projection,
as a non sequitur.

None of the communication failures in the novel
are failures of intelligence or good faith. They
are failures of projection compatibility.

\subsection{Repair-Induced Damage}

Every repair operator in the novel produces damage
alongside whatever repair it achieves. This follows
from the structure of the situation.

When the anomaly set of a repair operator omits
features of the system that are relevant to
stability, the operator will treat those features
as background conditions rather than as variables
its intervention will affect. When the intervention
proceeds, the omitted features respond to the
change in ways the operator did not anticipate.

The Fathers' repair of civilizational vulnerability
produces a population permanently dependent on a
control mechanism that prevents the development
of autonomous capacity. The repair addresses the
immediate anomaly while creating a structural
condition that becomes increasingly anomalous
over time.

The underground's repair of cognitive suppression,
were it to succeed at scale, would expose a
traumatized population to accurate information
without the institutional infrastructure required
to support the transition.

Maxim's repair of the broadcast mechanism produces
the withdrawal crisis, the political vacuum, the
economic disruption, and the invitation to imperial
invasion that the Wanderer had been working to
prevent.

The mutants' repair of their community's survival
produces a permanent withdrawal from any
participation in the broader situation, ensuring
that the most epistemically sophisticated
non-aligned actors remain unavailable for any
intervention that might otherwise benefit from
their analysis.

In each case the repair operator is locally
coherent. In each case its application produces
consequences that the operator's anomaly set
did not contain.

The political conflict on Saraksh is therefore
not a conflict between those who understand the
system and those who do not. It is a conflict
among repair programs each of which understands
part of the system well enough to intervene
effectively within that part, and none of which
understands the full system well enough to
predict the consequences of its own success.

This is the situation into which Rudolf Sikorski
has inserted himself, with more information,
longer time horizons, and more sophisticated
intervention theory than any other actor in
the novel.

His failure, examined in the chapter that follows,
is therefore not a personal failure.

It is a demonstration.

% ---------------------------------------------------------------
\section{Rudolf Sikorski and the Limits of Repair}
% ---------------------------------------------------------------

The figure of Rudolf Sikorski occupies an unusual
position in the architecture of \textit{Prisoners
of Power}.

He appears late. He operates in disguise. His
objectives remain opaque until the novel's final
pages. When he finally reveals himself and his
analysis, the narrative has already ended in the
sense that matters most: the Control Center has
been destroyed, the broadcasts have ceased, and
the consequences are already propagating through
the system in ways that cannot be reversed.

This timing is not a dramatic convenience. It is
the novel's central structural argument. The most
sophisticated actor in the system arrives too late
to prevent the least sophisticated actor's
intervention from determining the future.

\subsection{The Progressor as Theoretical Construct}

The Progressor represents the Noon Universe's
best answer to the intervention problem. Not a
conqueror, not a missionary, not a development
economist imposing external models, but a
sophisticated analyst operating from within the
system with a long time horizon and a commitment
to indigenous development rather than external
imposition.

Sikorski is among the most experienced of these
agents. His presence on Saraksh is not an accident
of the plot but a theoretical statement: the novel
is asking what happens when the best available
theory of intervention encounters a system of
sufficient complexity.

\subsection{What Sikorski Understands}

Sikorski understands the broadcast system at a
technical and functional level that no native
actor possesses. He understands the geopolitical
situation within a context invisible to anyone
operating from within the planet's own information
environment. He understands intervention theory.

He knows that rapid systemic change in a fragile
civilization produces consequences that outrun the
capacity of existing institutions to absorb them.
He knows that the removal of control mechanisms
without replacement stabilizing structures
generates the very crises that the mechanisms
were originally installed to prevent.

His strategy reflects this understanding. He does
not attempt to destroy the broadcasts, overthrow
the Fathers, or accelerate the underground's
revolutionary program. He works at the margins
of the system, creating conditions that will
eventually make the control architecture
unnecessary by building the institutional and
cognitive infrastructure that would allow the
population to function without it.

\subsection{The Structure of Sikorski's Failure}

Sikorski's plan is destroyed by Maxim.

The mechanism of destruction is not military,
political, or ideological. It is projective.

Maxim operates from a framework that does not
contain the distinctions Sikorski's analysis
requires. He cannot see the load-bearing function
of the broadcasts because his framework preserves
the distinction between freedom and oppression
while collapsing the distinction between constraint
and support. He cannot see Sikorski's long-term
developmental strategy because it is invisible
by design, operating below the threshold of
political visibility.

The collision between Maxim's framework and
Sikorski's framework is not a collision between
ignorance and knowledge. It is a collision between
two projections that preserve different distinctions
and therefore generate different anomaly sets and
different repair operators.

Neither projection is complete. The difference is
that Sikorski's projection is organized around
the consequences of interventions across long time
horizons, while Maxim's projection is organized
around the justice of the present state.

Both orientations track something real. Their
incompatibility is not resolvable by simply adding
one agent's information to the other's framework,
because the information each agent possesses is
structured by the projection that organized
its collection.

\subsection{Intervention, Control, and Irreversibility}

Irreversibility transforms the structure of the
intervention problem fundamentally.

In a reversible system, the cost of intervention
under incomplete observability is bounded by the
cost of correction. Agents can afford to act under
uncertainty because errors can be identified and
addressed.

In an irreversible system, the cost of intervention
under incomplete observability is potentially
unbounded. A single inadmissible intervention may
initiate trajectories that no subsequent action
can redirect toward the original target.

Sikorski understands this. His strategy of
incremental, invisible intervention is designed
precisely to preserve reversibility. Small changes
can be observed, evaluated, and adjusted. Large
rapid changes cannot.

Maxim does not understand this. His framework does
not preserve the distinction between reversible
and irreversible interventions as a primary
evaluative criterion.

\subsection{What Sikorski's Failure Means for Repair Theory}

The analysis of Sikorski's situation reveals a
limitation in repair theory as currently formulated.

Existing formulations assume that repair operators
can be designed independently and applied
sequentially. This model does not account for the
interaction between repair operators designed by
different agents operating from different
projections in the same system simultaneously.

On Saraksh, multiple repair programs are active
at the same time. Each program is designed to
address a real anomaly. Each program produces
consequences that affect the conditions under
which other programs operate. The programs are
not coordinated. They cannot be coordinated,
because the projections from which they are
designed do not preserve sufficient common
structure to support coordination.

The novel suggests that sufficiently complex
social systems require a theory of intervention
that accounts for the simultaneous activity of
multiple incompatible repair programs and the
interactions among them.

Current formal frameworks do not provide this theory.

The Wanderer understood the problem more clearly
than any other actor in the novel. He did not
solve it. Whether the problem is solvable in
principle is a question the novel poses without
answering. It is also a question that the
mathematics of admissibility and reachability
has only begun to approach.

% ---------------------------------------------------------------
% PART IV
% ---------------------------------------------------------------
\newpage
\part*{Part IV: The Ending and Its Implications}
\addcontentsline{toc}{part}{Part IV: The Ending and Its Implications}

% ---------------------------------------------------------------
\section{The Ending as Reachability Problem}
% ---------------------------------------------------------------

The final pages of \textit{Prisoners of Power}
refuse the satisfactions that the preceding
narrative has repeatedly made available and
then withdrawn.

Maxim has destroyed the Control Center. The
broadcasts have ceased. The Wanderer has delivered
his accounting of the consequences already in
motion. Famine, anarchy, institutional collapse,
the withdrawal crisis affecting perhaps a fifth
of the population, and the looming invasion from
the Island Empire whose suppression the Wanderer
had been planning through mechanisms now unavailable.

Maxim listens. He does not recant.

He acknowledges the consequences. He acknowledges
that he did not foresee them. He acknowledges
that the Wanderer's analysis is probably correct
in its broad outlines.

And then he refuses to leave.

He will stay on Saraksh. He will help stabilize
the situation. He remains glad that he destroyed
the Control Center because now the people can
be in charge of their own destiny.

The novel ends on this refusal. Not on resolution.
Not on vindication. Not on condemnation. On the
coexistence of two formally valid positions that
the narrative declines to adjudicate.

\subsection{Two Readings of the Same Event}

The destruction of the Control Center admits two
interpretations that the novel holds simultaneously
without resolving.

The first interpretation is Maxim's. The broadcast
system was a mechanism of neurological coercion
imposed on a population without consent. It
systematically prevented the population from
forming accurate representations of their own
situation. Whatever instabilities follow from its
removal are instabilities the population is now
at least capable of navigating as agents rather
than as managed objects.

The second interpretation is the Wanderer's. The
broadcast system was a load-bearing constraint in
a civilization that had recently nearly destroyed
itself. The instabilities following from its
removal are not navigable by a population that
has organized its psychological and social life
around the system's rhythms for generations. The
removal does not produce self-determination. It
produces disorientation at civilizational scale.

Both interpretations are accurate descriptions
of the same event. They do not contradict each
other at the level of fact. They diverge at the
level of which facts are assigned priority within
the evaluative framework each interpreter employs.

The formal analysis that follows does not resolve
this divergence. It explains why the divergence
is not resolvable from within the information
available to any actor in the system.

\subsection{Formal Reachability and Effective Reachability}

To make the disagreement precise, it is necessary
to distinguish two quantities that political
analysis typically conflates.

Let $\mathcal{A}_C$ denote the admissible region
of a civilization under a constraint field $C$,
and let $\mathcal{A}_0$ denote the admissible
region after removal of the constraint. Define
the formal reachability volume as
\[
V_R(\mathcal{A}) = \int_{\mathcal{A}} w(x)\, d\mu(x),
\]
where $w(x)$ is a viability weighting function
and $\mu$ is the induced measure on the admissible
manifold.

The naive interpretation of liberation assumes
$\mathcal{A}_C \subseteq \mathcal{A}_0$ and
therefore $V_R(\mathcal{A}_0) \geq
V_R(\mathcal{A}_C)$. Removing a constraint
expands the admissible region. The argument
appears to close here. It does not close here.

Let $N(x,t)$ denote the navigational capacity
of agents occupying state $x$ at time $t$.
Define the effective reachability volume as
\[
V_E(\mathcal{A},t) =
\int_{\mathcal{A}} N(x,t)\, w(x)\, d\mu.
\]

\begin{proposition}[Reachability Tradeoff]
\label{prop:tradeoff}
There exist systems for which
\[
V_R(\mathcal{A}_0) \geq V_R(\mathcal{A}_C)
\]
while simultaneously
\[
V_E(\mathcal{A}_0, t_0) < V_E(\mathcal{A}_C, t_0)
\]
for some interval $[t_0, t_1]$ immediately
following constraint removal.
\end{proposition}

\begin{proof}
Constraint removal enlarges the admissible
manifold, so formal reachability does not
decrease. However, navigational capacity evolves
on a slower timescale than constraint geometry.
Immediately after removal,
\[
N(x, t_0) \approx N_C(x),
\]
where $N_C$ was adapted to the previous constraint
structure. The enlarged manifold $\mathcal{A}_0$
contains regions for which $N(x, t_0) \approx 0$,
since agents have not yet developed the capacity
to navigate states that were previously inadmissible.
The contribution of these regions to effective
reachability remains negligible despite their
formal availability. Therefore formal reachability
increases while effective reachability decreases
over $[t_0, t_1]$:
\[
V_R \uparrow, \qquad V_E \downarrow. \qquad \square
\]
\end{proof}

Proposition~\ref{prop:tradeoff} formalizes the
disagreement between Maxim and the Wanderer with
precision that neither character possesses. Maxim
argues from formal reachability. The Wanderer
argues from effective reachability. They are not
disagreeing about the same quantity. They are
computing different things and calling both of
them liberation.

\subsection{The Lamphrodyne Interpretation}

Within RSVP field theory, the situation on Saraksh
after the broadcasts cease can be characterized
as a lamphrodyne event~\cite{flyxion2026rsvp}.

Let $\Phi(x,t)$ represent the local density of
cognitive possibility at position $x$ and time $t$,
and let $\mathbf{v}(x,t)$ represent the transport
field carrying navigational capacity through the
social manifold. The broadcast system operates
as a constraint on $\Phi$, suppressing possibility
density in regions of cognitive space corresponding
to accurate representations of the political system.
The transport field $\mathbf{v}$ has been calibrated
over generations to the constraint geometry.

\begin{theorem}[Lamphrodyne Overshoot]
\label{thm:lamphrodyne}
Let $\tau_C$ denote the timescale of constraint
removal and $\tau_v$ the adaptation timescale
of the transport field. If
\[
\tau_C \ll \tau_v,
\]
then there exists an interval $[t_0, t_1]$ during
which expansion of possibility exceeds transport
capacity:
\[
\frac{d\Phi}{dt} \gg \left|\nabla \cdot \mathbf{v}\right|.
\]
During this interval the system enters a state of
lamphrodyne overshoot.
\end{theorem}

\begin{proof}
Removal of the constraint produces a discontinuity
in the available possibility field,
\[
\Delta\Phi > 0,
\]
on timescale $\tau_C$. The transport field evolves
continuously on timescale $\tau_v$. Since
$\tau_C \ll \tau_v$,
\[
\mathbf{v}(t_0^+) \approx \mathbf{v}(t_0^-).
\]
Immediately after removal,
\[
\frac{d\Phi}{dt} \gg \frac{d\mathbf{v}}{dt}.
\]
The possibility field therefore expands faster
than navigational capacity can be transported
through it. Regions become formally reachable
before institutions can support navigation toward
them. The resulting mismatch constitutes
lamphrodyne overshoot.~$\square$
\end{proof}

Theorem~\ref{thm:lamphrodyne} characterizes the
Wanderer's prediction of mass withdrawal psychosis,
institutional collapse, and civilizational
disorientation as a prediction of lamphrodyne
overshoot. What the theorem provides is a
characterization precise enough to support the
following question: under what conditions does
lamphrodyne overshoot remain bounded, and under
what conditions does it propagate until it destroys
the navigational infrastructure of the system
entirely? That question the novel cannot answer.
It is an empirical question about the trajectory
of Saraksh after the final page, and also the
question that makes the novel's ending the
beginning of a research program rather than
a conclusion.

\subsection{Constraint Removal Non-Monotonicity}

Most political theories rest on an implicit
assumption so obvious it rarely receives
explicit statement:
\[
C_1 < C_2 \;\Longrightarrow\; V_R(C_1) > V_R(C_2).
\]
Fewer constraints imply greater freedom. Greater
freedom implies greater reachability. The direction
of the inequality is taken as self-evident.

The ending of \textit{Prisoners of Power} directly
challenges this assumption, and the challenge can
be stated with mathematical precision.

Let $V_R(C)$ denote the formal reachability volume
of a civilization under constraint intensity $C$,
and let $S(C)$ denote the stabilizing infrastructure
available to the population at that constraint level.
Define effective reachability as the product
\[
V_E(C) = V_R(C)\cdot S(C).
\]

The standard liberation assumption treats $S$ as
independent of $C$: removing a constraint expands
$V_R$ without affecting $S$. Under this assumption,
$V_E$ is monotonically decreasing in $C$ and the
Wanderer's position is simply wrong.

The novel proposes a different relationship.

\begin{theorem}[Constraint Removal Non-Monotonicity]
\label{thm:nonmonotone}
Suppose stabilizing infrastructure is itself
partially produced by the constraint field, so that
\[
\frac{dS}{dC} > 0.
\]
Then $V_E(C) = V_R(C)\cdot S(C)$ need not be
monotonically decreasing in $C$. In particular,
there may exist a critical value $C^\star$ satisfying
\[
\frac{dV_E}{dC}\bigg|_{C^\star} = 0,
\]
at which effective reachability is locally maximal.
Consequently,
\[
V_E(C^\star) > V_E(0).
\]
A partially constrained civilization can possess
strictly greater effective reachability than a
completely unconstrained one.
\end{theorem}

\begin{proof}
Write $V_E(C) = V_R(C)\cdot S(C)$.
Differentiating with respect to $C$,
\[
\frac{dV_E}{dC}
= \frac{dV_R}{dC}\cdot S(C)
+ V_R(C)\cdot \frac{dS}{dC}.
\]
Since increasing constraint reduces formal
reachability, $dV_R/dC < 0$.
By hypothesis, $dS/dC > 0$.
The two terms therefore have opposite sign.

At $C = 0$, the constraint is absent and
$S(0)$ may be small or zero if stabilizing
infrastructure depends on the constraint for
its production. At large $C$, $V_R(C)$ is
severely suppressed. Between these extremes,
the product $V_R(C)\cdot S(C)$ attains an
interior maximum at $C^\star$ satisfying
$dV_E/dC = 0$, i.e.
\[
\frac{dV_R}{dC}\cdot S(C^\star)
= -V_R(C^\star)\cdot \frac{dS}{dC}\bigg|_{C^\star}.
\]
At this point $V_E(C^\star) > V_E(0)$
whenever $S(0)$ is sufficiently small relative
to $S(C^\star)$.~$\square$
\end{proof}

Theorem~\ref{thm:nonmonotone} does not establish
that the Wanderer is correct. It establishes
something more important: that Maxim's position
is not mathematically trivial. If the broadcast
system has been generating stabilizing
infrastructure over the decades of its operation
--- if the cognitive habits, institutional
routines, and social coordination mechanisms
of Saraksh's population have formed around
the constraint rather than independently of it
--- then the assumption $dS/dC > 0$ holds,
and the constraint removal non-monotonicity
theorem applies.

Whether $C^\star$ is close to the actual
broadcast intensity, and whether Sikorski's
incremental program was designed to reduce $C$
toward $C^\star$ rather than toward zero, are
empirical questions the novel leaves open.

What the theorem settles is the logical structure
of the disagreement: both positions are
consistent with the mathematics.

\subsection{The Admissibility--Stability Tradeoff}

The preceding analysis treats the disagreement
between Maxim and the Wanderer as a disagreement
about facts: what the broadcasts actually do,
what their removal would actually produce. But
there is a deeper disagreement running beneath
the factual one, and it is of a different kind.

Throughout the novel, two distinct optimization
targets are in play.

Let $A(t)$ denote epistemic admissibility: the
degree to which the population's representational
frameworks preserve the distinctions required
for accurate understanding of their situation.
Let $P(t)$ denote population stability: the
degree to which the civilization maintains
coherent institutional function, physical
welfare, and social coordination.

The Unknown Fathers maximize $P(t)$ at the
cost of $A(t)$. The broadcasts suppress
admissibility in order to maintain stability.

Maxim maximizes $A(t)$ at the cost of $P(t)$.
He destroys the mechanism of admissibility
suppression even at the cost of the stability
it was maintaining.

Suppose a civilization's welfare can be expressed
as a utility functional of the form
\[
U = \alpha A + \beta P,
\]
where $\alpha, \beta \geq 0$ are weighting
coefficients reflecting how the civilization
values epistemic admissibility relative to
population stability.

\begin{proposition}[Admissibility--Stability Tradeoff]
\label{prop:tradeoff2}
The disagreement between Maxim and the Wanderer
is equivalent to a disagreement about the
coefficients $(\alpha, \beta)$ in the utility
functional $U = \alpha A + \beta P$. In
particular, no accumulation of shared factual
information about the state of Saraksh can
resolve a disagreement that is constituted
by different utility geometries.
\end{proposition}

\begin{proof}
Suppose Maxim and the Wanderer share complete
information about the current state of Saraksh
and agree on the causal consequences of
all available interventions. Their disagreement
about what to do therefore cannot arise from
different beliefs about facts or consequences.

If they nevertheless recommend different
interventions, this must reflect different
evaluations of the same consequence profile.
Under the utility functional $U = \alpha A
+ \beta P$, Maxim's preference for epistemic
transparency implies a high ratio $\alpha/\beta$,
while the Wanderer's preference for civilizational
stability implies a low ratio $\alpha/\beta$.

For any finite body of factual information
$\mathcal{F}$, if both agents update
their beliefs identically on $\mathcal{F}$,
their posterior recommendations will still
diverge whenever $\alpha_{\text{Maxim}}/\beta_{\text{Maxim}}
\neq \alpha_{\text{Wanderer}}/\beta_{\text{Wanderer}}$.

Additional facts therefore cannot close the gap.
The conflict is not epistemic. It is
variational.~$\square$
\end{proof}

This result illuminates one of the novel's
most frustrating structural features: the
failure of argument between the factions.

Throughout \textit{Prisoners of Power}, characters
who share information do not converge in their
recommendations. Readers accustomed to treating
political disagreement as a consequence of
factual ignorance may find this inexplicable.
Proposition~\ref{prop:tradeoff2} provides the
explanation. When disagreement is variational
rather than epistemic, more information does
not resolve it. It merely sharpens the
contours of the divergence.

The Wanderer's exhaustion in the final scene
is not the exhaustion of someone who has been
proved wrong. It is the exhaustion of someone
who has been unable to communicate that the
other party is optimizing a different objective
function. That communication failure is not
rhetorical. It is mathematical.

\subsection{The Projection Horizon Theorem}

The three results developed above characterize
the structure of the disagreement at the novel's
ending. They do not yet explain why the novel
refuses to adjudicate between the positions.

The final result of this section provides
that explanation.

\begin{theorem}[Projection Horizon]
\label{thm:horizon}
Let $H$ denote the prediction horizon available
to an observer within the system, and let
$\Delta R(T)$ denote the difference in cumulative
effective reachability between two interventions
over evaluation horizon $T$. If $T > H$, then
the sign of $\Delta R(T)$ cannot generally be
determined from information available to any
observer within the system.
\end{theorem}

\begin{proof}
Evaluation of $\Delta R(T)$ requires knowledge
of both intervention trajectories across the
full horizon $T$. An observer with prediction
horizon $H < T$ has direct access only to
trajectory segments of length at most $H$.

Any estimate of $\Delta R(T)$ from horizon-$H$
information therefore requires a projection
operator $\pi_H$ mapping observable segments
to full-horizon trajectory estimates. Different
projection operators yield different estimates:
\[
\pi_H^{(1)}(\Delta R(T))
\;\neq\;
\pi_H^{(2)}(\Delta R(T))
\]
in general. Since no observer has access to
the unrealized counterfactual trajectory or
to exogenous disturbances $\eta(t)$ for
$t > H$, no canonical choice of $\pi_H$
is available.

Therefore $\operatorname{sign}(\Delta R(T))$
cannot be determined from within the system
for $T > H$.~$\square$
\end{proof}

Theorem~\ref{thm:horizon} captures the deepest
point of the novel's ending, and it is worth
stating carefully.

The problem is not that Maxim lacks information
that the Wanderer possesses. The problem is
that the relevant information does not yet
exist. It resides in the future. The quantity
required to adjudicate their disagreement ---
the sign of the difference in cumulative
effective reachability between the world where
Maxim acts and the world where he does not ---
lies beyond the prediction horizon of any
actor inside the system at the moment of
decision.

This is why the novel's unresolved ending
feels structurally necessary rather than
evasive. The Strugatskys have not withheld
an answer. They have constructed a situation
in which the answer is not yet available to
anyone inside the story.

The uncertainty is not psychological.
It is not a matter of insufficient research
or insufficient courage.
It is a structural feature of the problem:
the horizon of evaluation exceeds the horizon
of observation, and no projection operator
bridges the gap without importing assumptions
that are themselves contested.

Maxim knows this, at some level, when he
refuses to leave at the novel's end. He
cannot prove he was right. He also cannot
be proved wrong. He stays because the
question remains open, and staying is the
only form of commitment available to someone
who cannot compute the answer.

\subsection{The Counterfactual Indeterminacy Theorem}

The preceding results establish that constraint
removal can decrease effective reachability in
the short term (Proposition~\ref{prop:tradeoff}),
that rapid removal produces lamphrodyne overshoot
(Theorem~\ref{thm:lamphrodyne}), that the
liberation assumption is not mathematically
trivial (Theorem~\ref{thm:nonmonotone}), that
the disagreement between Maxim and the Wanderer
is variational rather than epistemic
(Proposition~\ref{prop:tradeoff2}), and that
the sign of the long-run reachability difference
lies beyond any actor's projection horizon
(Theorem~\ref{thm:horizon}). They do not
establish that Maxim's intervention was wrong.

To establish that, one would need to compare the
long-term trajectory generated by the intervention
against the long-term trajectory that would have
obtained without it. This comparison is not available.

\begin{theorem}[Counterfactual Indeterminacy]
\label{thm:indeterminacy}
No actor within the system possesses sufficient
information to determine the sign of
\[
\Delta R(T) = R(\tau_M, T) - R(\tau_W, T),
\]
where $\tau_M$ is the trajectory generated by
destruction of the Control Center, $\tau_W$ is
the trajectory generated by continuation of the
Wanderer's incremental program, and cumulative
effective reachability over horizon $T$ is defined
by
\[
R(\tau, T) = \int_0^T V_E(\tau(t), t)\, dt.
\]
\end{theorem}

\begin{proof}
Evaluation of $\Delta R(T)$ requires knowledge
of both trajectories across the full horizon $T$.
Only $\tau_M$ is realized. The counterfactual
trajectory $\tau_W$ is unobservable by definition.
Furthermore, both trajectories depend upon future
contingent events,
\[
\tau(t) = F(x_0, \eta(t)),
\]
where $\eta(t)$ represents exogenous disturbances
including military action, institutional adaptation,
ecological and material conditions, elite decisions,
and technological developments on timescales
ranging from months to generations. Since neither
Maxim nor the Wanderer has access to future values
of $\eta(t)$, and since the counterfactual $\tau_W$
is by definition not instantiated, neither can
compute $\Delta R(T)$. The quantity required to
adjudicate the dispute is therefore not merely
unknown but epistemically inaccessible to any
actor operating within the system.~$\square$
\end{proof}

Theorem~\ref{thm:indeterminacy} establishes that
the disagreement between Maxim and the Wanderer
is not a disagreement that more information would
resolve. Additional information about the current
state of Saraksh would not provide access to the
counterfactual trajectory.

The novel ends without resolving a question that
is formally unresolvable. This is not irresolution.
It is precision. The reader who finishes the novel
wanting a verdict has misidentified the question.

The question is not what should have been done.

The five results of this section jointly explain
why no actor within the system could answer it.
Constraint removal is not monotone in effective
reachability (Theorem~\ref{thm:nonmonotone}),
so the direction of the liberation argument is
not self-evident. The disagreement between Maxim
and the Wanderer is variational rather than
epistemic (Proposition~\ref{prop:tradeoff2}),
so shared facts cannot resolve it. The sign of
the long-run reachability difference lies beyond
any actor's prediction horizon
(Theorem~\ref{thm:horizon}), so the relevant
quantity is not computable from within the system.
Rapid constraint removal produces lamphrodyne
overshoot (Theorem~\ref{thm:lamphrodyne}), so
short-term evidence is not diagnostic of
long-run outcomes. And cumulative effective
reachability across both trajectories is
epistemically inaccessible
(Theorem~\ref{thm:indeterminacy}), so even
retrospective evaluation remains underdetermined.

The Strugatskys did not construct an ending
that refuses resolution because they could not
decide. They constructed an ending that refuses
resolution because the question it poses is
structurally undecidable from within the
information available to any actor at the moment
of decision.

That is the research program the following
chapter begins to sketch.

% ---------------------------------------------------------------
% PART V
% ---------------------------------------------------------------
\newpage
\part*{Part V: Beyond Admissibility}
\addcontentsline{toc}{part}{Part V: Beyond Admissibility}

% ---------------------------------------------------------------
\section{Beyond Admissibility}
% ---------------------------------------------------------------

The preceding chapters have developed a reading
of \textit{Prisoners of Power} through the concepts
of projection, admissibility, reachability, and
repair. This reading has produced genuine results.
The novel's recurring pattern of model collapse
is more precisely characterized as projection
failure than as ignorance or naivety. The broadcast
system is more precisely characterized as an
admissibility architecture than as propaganda.
The political conflict is more precisely
characterized as a collision among incompatible
repair programs. The ending is more precisely
characterized as a reachability problem than as
a political resolution.

Yet the preceding analysis has also repeatedly
encountered phenomena that the framework handles
awkwardly or does not handle at all. The purpose
of this chapter is to examine those phenomena
directly. Not to explain them away. Not to extend
the framework prematurely to cover them. But to
characterize as precisely as possible the boundary
between what the current mathematics can reach
and what lies beyond it.

That boundary is itself a theoretical result.

\subsection{Loyalty}

The first phenomenon that resists formalization
is loyalty.

Throughout the novel, characters act in ways that
cannot be fully explained by their assessment of
outcomes. Maxim's companions in the underground
maintain their commitments under conditions where
rational reachability analysis would recommend
withdrawal. Soldiers follow orders whose
consequences they can observe to be catastrophic.
Members of the resistance accept imprisonment and
death for objectives whose achievement they are
unlikely to witness.

One might attempt to absorb loyalty into the
framework by expanding the objective function:
the agent values consistency of commitment
independently of outcomes. This move is formally
possible. It is also explanatorily empty. To say
that an agent values loyalty is to redescribe
the phenomenon without explaining it. The question
is why loyalty functions as a stable attractor in
human social systems, why it resists revision
under conditions where outcome-based reasoning
would recommend revision, and what role it plays
in the stability and navigability of the social
systems within which it appears.

These questions require a theory of motivational
structure that the current framework does not provide.

\subsection{Charisma}

The second phenomenon is charisma.

Maxim himself is an example. He repeatedly
persuades individuals to assist him, join him,
or trust him in situations where rational
assessment of his track record would not support
such trust. His persuasiveness is not primarily
a function of the arguments he makes. It is a
function of something the text conveys as personal
presence, energy, and conviction.

Charisma of this kind has observable effects on
the reachability volumes of social systems. Leaders
who possess it can mobilize collective action that
would not occur in their absence, sustain
commitment through periods of failure, and produce
coordination among agents whose individual
assessments would not generate coordinated action.
Yet charisma does not appear naturally within the
formal vocabulary developed in this essay. It is
not a property of a state, a transition, an
anomaly set, or a repair operator. It is a property
of the relationship between a particular agent and
the social field within which that agent operates.

Whether this relationship can be formalized within
a geometric framework remains an open question.

\subsection{Fear}

The third phenomenon is fear.

Fear in \textit{Prisoners of Power} is not simply
a negative utility assigned to certain outcomes.
It is a cognitive and motivational state that
restructures the space of actions an agent is
capable of considering. Under sufficient fear,
options that are formally available become
practically inaccessible.

This contraction is not identical to the
contraction produced by broadcast constraint.
Broadcast constraint operates externally, shaping
the representational space through which the agent
processes information. Fear operates internally,
altering the relationship between the agent's
representations and their motivational structure.

Part of what the Wanderer is predicting when he
describes the consequences of Maxim's action is
not merely institutional collapse but psychological
collapse: a population that has organized its
affective life around the management of fear,
and whose management strategies have been suddenly
rendered unavailable without replacement. The
formal framework can identify that practical
navigability and formal reachability have diverged.
It cannot yet characterize the internal structure
of the divergence in terms that would support the
design of interventions specifically targeted at
restoring navigability.

\subsection{Historical Momentum}

The fourth phenomenon is historical momentum.

At several points in the novel, events proceed in
ways that no actor intended and that the preferences
and actions of individual actors do not fully
explain. The war against Honti is initiated through
a process in which each participant follows a
locally rational path toward a collectively
irrational outcome. No single actor decides that
the war should occur. The war occurs anyway because
the interaction of multiple repair programs, each
responding to real anomalies within its own
framework, produces a trajectory that none of the
programs was designed to generate.

Reachability theory can describe the phenomenon
in retrospect. What it does not yet provide is
a systematic method for identifying, in advance,
which interactions among simultaneously active
repair programs will produce emergent trajectories
that none of the programs was designed to generate.

Historical momentum is not noise. It is the
structured interaction of multiple repair programs
operating from incompatible projections in a system
where the interaction terms are as causally
significant as the individual programs themselves.
A theory adequate to the phenomenon would need to
characterize the interaction structure, not merely
the individual programs.

\subsection{Moral Commitment}

The fifth phenomenon is moral commitment.

Maxim's refusal to recant at the novel's end is
not primarily a reachability judgment. He does not
say that he has calculated superior outcomes. He
says that the people can now be in charge of their
own destiny, and that this matters regardless of
the immediate consequences. This is a statement
about the intrinsic value of self-determination
that is not reducible to its effects on reachable
futures.

One could attempt to treat moral commitment as
a constraint on the objective function. This
treatment is formally coherent but leaves the
most important question unanswered. Where do
the constraints come from? Why do they persist
under pressure? Why do some moral commitments
function as stable attractors while others
dissolve under adversity? And why do individuals
sometimes sacrifice their own reachable futures
in service of moral commitments that benefit
others they will never meet?

\subsection{Narrative Agency}

The sixth phenomenon is narrative agency.

Throughout the novel, characters act not only in
response to their assessment of the current state
but in response to their understanding of the
story they are in. Maxim repeatedly makes decisions
that are influenced by the kind of protagonist he
believes himself to be. The Robinson Crusoe frame
is not merely an interpretive error. It is a
motivational structure that organizes which actions
feel available, which feel appropriate, and which
feel narratively coherent.

Agents do not simply optimize over reachable
futures. They occupy narrative positions that shape
which futures they are capable of orienting toward.

Whether narrative can be incorporated into a
geometric theory of reachability, and if so through
what formal machinery, is among the most
interesting open questions the novel raises for
the framework.

\subsection{The Boundary and Its Significance}

The six phenomena examined in this chapter share
a common structure. Each is causally relevant to
the trajectories of social systems. Each influences
which futures become reachable and which do not.
Each appears in the novel as a force that shapes
outcomes in ways that purely structural analysis
of state spaces and transition operators does not
fully capture. And each resists clean formalization
within the current framework.

This resistance is not a reason to abandon the
framework. It is a reason to treat the framework
as a first approximation rather than a complete
theory.

The boundary between what the mathematics can
reach and what lies beyond it is not a defect.
It is a map of where the mathematics needs to grow.

The Strugatskys did not possess formal theories
of admissibility, reachability, or repair. They
possessed something else: the ability to construct
imagined worlds complex enough that all of these
phenomena appear together, interact with one
another, and resist reduction to any single
explanatory frame.

The novel sees further than the theory.

That is not a failure of the theory. It is the
most useful thing a literary work can do for a
theoretical program: identify the terrain that
the program has not yet reached.

% ---------------------------------------------------------------
% PART VI
% ---------------------------------------------------------------
\newpage
\part*{Part VI: Literature as Theory Discovery}
\addcontentsline{toc}{part}{Part VI: Literature as Theory Discovery}

% ---------------------------------------------------------------
\section{Literature as Theory Discovery}
% ---------------------------------------------------------------

This essay began with a methodological claim.

Literary works sometimes function not as
illustrations of existing theories but as
instruments for discovering phenomena that
existing theories struggle to describe. The value
of such works for theoretical inquiry lies not
in the examples they provide but in the resistance
they offer: the places where the narrative sees
further than the formalism, where the imagined
world contains structure that the available
mathematics cannot yet reach.

\textit{Prisoners of Power} has proved to be
such a work.

The novel's treatment of projection failure is
more precise than most formal accounts of epistemic
error. The Strugatskys understood, without the
vocabulary of admissibility theory, that the
primary site of cognitive failure is not the
reasoning process but the representational
framework within which reasoning occurs. Maxim
does not fail because he reasons badly. He fails
because the categories through which he reasons
do not preserve the distinctions the world
requires.

The novel's treatment of constraint removal is
more nuanced than most formal accounts of
liberation. The Strugatskys understood, without
the vocabulary of lamphrodyne dynamics~\cite{flyxion2026rsvp},
that the removal of a load-bearing constraint
produces a different situation than the removal
of a merely restrictive constraint. They
understood that formal expansion of reachability
and practical expansion of navigability are
different quantities that can move in opposite
directions (Proposition~\ref{prop:tradeoff}).
They understood that the liberation assumption
--- fewer constraints imply greater freedom ---
fails when stabilizing infrastructure depends
on the constraint for its production
(Theorem~\ref{thm:nonmonotone}). And they
understood that the disagreement between Maxim
and the Wanderer cannot be resolved by sharing
more information because it is constituted by
different utility geometries rather than
different factual beliefs
(Proposition~\ref{prop:tradeoff2}).

The novel's treatment of intervention under
incomplete observability is more rigorous than
most formal accounts of political action. The
Strugatskys understood, without the vocabulary
of counterfactual indeterminacy, that the question
of whether an intervention was correct cannot be
answered from within the realized trajectory.
They understood that the most sophisticated actor
in a complex system cannot guarantee admissible
outcomes because the interaction of multiple
incompatible repair programs operating
simultaneously produces emergent trajectories
that no single program was designed to generate.

These understandings were not expressed as
theorems. They were expressed as narrative
structure, character behavior, and the specific
shape of the novel's refusal to resolve its
central questions.

\subsection{Two Modes of Inquiry}

The relationship between narrative and mathematics
is not one of translation.

Mathematics achieves precision by abstracting
away from particularity. A theorem about
reachability tradeoffs applies to any system
satisfying the relevant conditions, regardless
of whether that system involves Soviet-era
broadcast technology, medieval religious
institutions, contemporary social media
architectures, or colonial governance structures.
The abstraction is the source of the generality.
But the abstraction is also a loss.

The formal framework developed in this essay
cannot represent Maxim Kammerer as an individual.
It can represent a projection operator, an anomaly
set, a repair program. It cannot represent the
specific texture of a young man from an idealistic
civilization encountering, for the first time,
a world in which good intentions produce
catastrophic consequences.

That texture is not decoration. It is information.

The specific way in which Maxim's confidence
collapses, the specific way in which he nevertheless
refuses to recant, the specific way in which the
Wanderer's exhaustion reads as the exhaustion of
someone who has understood the problem for longer
than Maxim has been alive: these particulars carry
theoretical content that the formal framework
cannot capture without losing what makes the
content significant.

Literature achieves depth by preserving
particularity. The novel can hold simultaneously,
without resolving, the validity of Maxim's position
and the validity of the Wanderer's position,
because narrative can sustain contradiction in
a way that formal proof cannot. The ending of
\textit{Prisoners of Power} is not a failure of
resolution. It is the representation of a genuine
contradiction whose resolution would require
information that is structurally inaccessible.

The two modes are therefore complementary rather
than competitive. Mathematics tells us what the
structure of the problem is. Narrative tells us
what it feels like to inhabit the problem. Neither
is sufficient alone. Together they illuminate more
than either can reach independently.

\subsection{What the Strugatskys Discovered}

The claim that a literary work can constitute
theoretical discovery requires a standard of
evidence. It is not sufficient to observe that
a novel contains themes that happen to resemble
concepts in a formal theory.

The stronger claim requires demonstrating that
the novel exhibits structural features that cannot
be fully accounted for by the theoretical
frameworks available at the time of its composition,
and that the specific structure of these features
anticipates distinctions that subsequent
theoretical work has found it necessary to introduce.

\textit{Prisoners of Power} satisfies this standard
in at least three respects.

The distinction between formal and effective
reachability, which Proposition~\ref{prop:tradeoff}
formalizes, is not present in the political theory
or cognitive science literature of 1969. The novel
embeds this distinction in its plot structure: the
entire tragedy of the ending depends on the gap
between what the population can now formally access
and what they can actually navigate toward. The
distinction is not stated. It is enacted,
repeatedly, at the level of narrative consequence.

The distinction between constraint removal and
load-bearing constraint, which the lamphrodyne
analysis formalizes, is present in systems theory
literature of the period but not in the political
theory literature that would have been the obvious
reference for a novel about totalitarianism. The
Strugatskys embed this distinction in the
Wanderer's argument without reducing it to the
political categories their readers would have
found familiar.

The structure of counterfactual indeterminacy,
which Theorem~\ref{thm:indeterminacy} formalizes,
is present in philosophy of science literature
but not in the narrative conventions of political
fiction. Political fiction conventionally resolves
by revealing which side was correct. The
Strugatskys construct a situation in which this
question is not merely unresolved but unresolvable,
and they construct it precisely enough that the
unresolvability feels like a demonstration
rather than an evasion.

\subsection{The Critical Reception as Evidence}

The history of the novel's critical reception
provides indirect evidence for this claim.

Critical interpretation of \textit{Prisoners of
Power} has consistently proceeded by locating the
novel within existing conceptual frameworks:
dystopian fiction, Soviet allegory, Cold War
anxiety, outsider intervention narrative,
anti-totalitarian critique. These classifications
are not incorrect. But the critical literature
also consistently reports a residue: something
that the available frameworks do not quite capture,
a quality of unresolvedness that readers experience
as significant rather than merely incomplete.

This residue is evidence that the novel contains
structure that existing frameworks consistently
omit. The frameworks locate what is familiar.
The residue marks what is not.

The analysis developed in this essay has argued
that the residue is not a vague aesthetic quality
but a specific theoretical content: the novel's
precise treatment of projection failure,
load-bearing constraint, and counterfactual
indeterminacy in the context of intervention
under incomplete observability.

Critical readers have sensed this content without
possessing the vocabulary to articulate it. The
task of the present essay has been to develop
that vocabulary by placing the novel in
conversation with formal frameworks capable of
reaching the structure the critical tradition has
consistently felt but consistently failed to name.

\subsection{The Research Program}

The analysis has generated several open questions
that constitute a research program rather than
a conclusion.

The Reachability Tradeoff Proposition establishes
that formal and effective reachability can diverge.
It does not characterize the conditions under
which the divergence is temporary and
self-correcting versus the conditions under which
it initiates a cascade that permanently reduces
effective reachability. That characterization would
require a theory of navigational capacity
development: how do agents and institutions develop
the capacity to navigate newly available states,
at what rate, and under what conditions does the
development process fail?

The Lamphrodyne Overshoot Theorem establishes
that rapid constraint removal produces overshoot.
It does not characterize the conditions under
which the overshoot remains bounded versus the
conditions under which it propagates to systemic
collapse. That characterization would require a
theory of transport field resilience.

The Counterfactual Indeterminacy Theorem
establishes that the comparison between realized
and counterfactual trajectories is epistemically
inaccessible. It does not characterize what this
inaccessibility implies for intervention design.
If the consequences of intervention cannot be
compared to the consequences of non-intervention,
what standards should govern the decision to act?

These questions extend beyond the analysis of
\textit{Prisoners of Power}. They are questions
about intervention in complex social systems under
conditions of incomplete observability,
irreversibility, and multiple simultaneously active
repair programs. They are relevant to foreign
intervention, development policy, institutional
reform, ecological governance, and the alignment
of artificial systems with human values.

The Strugatskys wrote a novel about a young man
stranded on a damaged planet who destroys a mind
control system and produces consequences he did
not anticipate. They also wrote, without intending
to and possibly without knowing it, one of the
most precise fictional treatments of the geometry
of intervention under incomplete observability
that the twentieth century produced.

Whether that geometry can be fully formalized
remains to be seen. What this essay has attempted
to demonstrate is that the attempt to formalize
it is worthwhile, that the novel is a more precise
guide to the relevant terrain than the political
allegory reading suggests, and that the
conversation between narrative and mathematics,
when both are allowed to challenge one another,
produces insights that neither can reach alone.

The Strugatskys did not solve the intervention
problem. Neither has this essay. What the novel
did, and what this essay has attempted to do in
response, is make the problem visible in its full
complexity: not as a failure of political will or
moral clarity, but as a genuine difficulty in the
geometry of knowing and acting in worlds whose
most important structures cannot be directly
observed.

That visibility is the beginning of understanding.
It is also, for now, where understanding stops.

% ---------------------------------------------------------------
\section*{Conclusion: The Geometry of Control}
\addcontentsline{toc}{section}{Conclusion: The Geometry of Control}
% ---------------------------------------------------------------

\textit{Prisoners of Power} is not primarily a
novel about totalitarianism.

It is a novel about the conditions under which
any agent, operating within any system whose most
important structures remain partially hidden, can
act in ways that reliably improve rather than
damage the futures available to the inhabitants
of that system.

The answer the novel provides is not encouraging.

Every actor fails. The most sophisticated actor
fails last and most painfully, because his
understanding is sufficient to anticipate the
consequences of others' interventions without
being sufficient to prevent them.

The formal analysis developed in this essay has
not improved on this answer. It has, however,
made the structure of the failure more precise.

Maxim and the Wanderer are not disagreeing about
facts or values. They are computing different
quantities: formal reachability and effective
reachability, present constraint and future
navigability, realized trajectory and counterfactual
trajectory. The novel's refusal to adjudicate
between them is not irresolution. It is the
recognition that the quantity required to
adjudicate---cumulative effective reachability
across the relevant long-term horizon, compared
across realized and counterfactual
trajectories---is not accessible to any actor
operating within the system.

The geometry of control is therefore not the
geometry of power alone. It is the geometry of
knowledge, representation, and the relationship
between what can be observed and what can be
reached.

Saraksh is a world in which the atmosphere itself
is a projection operator, in which the political
system runs on a preference-field architecture,
in which every repair program generates anomalies
it was not designed to see, and in which the most
careful long-term intervention is destroyed by
a well-intentioned agent operating from an
incompatible and less adequate projection.

It is also, in these respects, a reasonably
accurate description of every sufficiently complex
social system that has ever existed.

The Strugatskys set their novel on an alien planet.

The planet is not alien.

% ---------------------------------------------------------------
% APPENDICES
% ---------------------------------------------------------------
\newpage
\appendix

\section{A Formal Sketch of Projection Failure}

This appendix collects and extends the formal
definitions introduced in the body of the essay.

\begin{definition}[State Space]
Let $\Omega$ be a measurable space representing
all possible states of a social system. Elements
$\omega \in \Omega$ include institutional
configurations, population belief distributions,
material conditions, and latent causal variables
not directly observable by any agent.
\end{definition}

\begin{definition}[Projection Operator]
A projection operator is a measurable map
$\pi: \Omega \to \Sigma$ from the full state
space to an agent's representational space
$\Sigma$. The fiber $\pi^{-1}(\sigma)$ for
$\sigma \in \Sigma$ represents the set of full
states that are indistinguishable to the agent.
\end{definition}

\begin{definition}[Decision Class]
A decision class $\mathcal{D}$ is a set of
functions $d: \Omega \to \mathcal{A}$ mapping
states to actions. A decision is correctly
implemented if the action $d(\omega)$ is chosen
at state $\omega$.
\end{definition}

\begin{definition}[Admissibility]
A projection $\pi$ is admissible for decision
class $\mathcal{D}$ if for every $d \in \mathcal{D}$
and every pair $\omega_1, \omega_2 \in \Omega$
such that $d(\omega_1) \neq d(\omega_2)$, we
have $\pi(\omega_1) \neq \pi(\omega_2)$.
Equivalently, $\pi$ is admissible for $\mathcal{D}$
if it does not identify any pair of states that
$\mathcal{D}$ requires to be distinguished.
\end{definition}

\begin{definition}[Projection Collapse]
A projection $\pi$ exhibits collapse with respect
to decision class $\mathcal{D}$ at state
$\omega$ if there exists $d \in \mathcal{D}$
and $\omega' \in \Omega$ such that
$\pi(\omega) = \pi(\omega')$ but
$d(\omega) \neq d(\omega')$.
\end{definition}

\begin{remark}
Projection collapse is the formal counterpart of
what the body of the essay calls projection
confidence failure. An agent suffering projection
collapse will behave as if states $\omega$ and
$\omega'$ are equivalent when they are not, and
will therefore select actions appropriate to one
when the situation calls for the other.
\end{remark}

\section{Reachability Partitions and Social Topology}

\begin{definition}[Reachability Relation]
Given a social system with state space $\mathcal{S}$
and action repertoire $\mathcal{I}$, define the
reachability relation $\sim_T$ by: $s \sim_T s'$
if and only if there exists a sequence of
interventions $\iota_1, \ldots, \iota_k \in
\mathcal{I}$ such that $\iota_k \circ \cdots
\circ \iota_1(s)$ places positive probability
on a neighborhood of $s'$ within time horizon $T$.
\end{definition}

\begin{definition}[Reachability Partition]
The reachability partition of $\mathcal{S}$ under
$\sim_T$ is the partition of $\mathcal{S}$ into
equivalence classes. Each class is a reachability
component: the set of states mutually accessible
from one another within horizon $T$.
\end{definition}

\begin{remark}
Different agents acting from different projections
effectively inhabit different reachability
components even when occupying the same physical
state. An agent whose projection collapses
distinctions necessary to navigate toward a target
state cannot reach that state regardless of
whether it is formally reachable. The effective
reachability component of an agent is therefore
the intersection of the formal reachability
component with the set of states the agent's
projection permits them to navigate toward.
\end{remark}

The factions of Saraksh---the Unknown Fathers,
the underground, the mutants, the military, the
Wanderer---inhabit effectively different
reachability components despite occupying the
same physical territory. Their political conflicts
can be understood as disputes about which
reachability component should be expanded,
contracted, or dissolved.

\section{Repair Operators on Political Systems}

\begin{definition}[Repair Operator]
A repair operator is a map $\rho: \mathcal{S}
\to \Delta(\mathcal{S})$ from states to
distributions over states. A repair operator
is designed to reduce the measure of some
anomaly set $\mathcal{A} \subseteq \mathcal{S}$
as assessed by the agent who designed it.
\end{definition}

\begin{definition}[Repair Compatibility]
Two repair operators $\rho_1$ and $\rho_2$
designed by agents with projections $\pi_1$
and $\pi_2$ are compatible if the anomaly sets
$\mathcal{A}_1$ and $\mathcal{A}_2$ do not
contain states that the opposing operator's
repair trajectory passes through.
\end{definition}

\begin{remark}
The repair operators of the major factions in
\textit{Prisoners of Power} are pairwise
incompatible in this sense. Each operator's
repair trajectory passes through states the
other operators identify as anomalies. The
resulting dynamic is not convergent. Each
repair induces new anomalies that motivate
further repair by other operators, producing
the oscillatory instability that characterizes
the novel's political landscape.
\end{remark}

\begin{definition}[Repair-Induced Anomaly]
Given repair operator $\rho$ targeting anomaly
set $\mathcal{A}$ from projection $\pi$,
define the repair-induced anomaly set as
\[
\mathcal{A}^\rho = \{s' \in \mathcal{S} :
s' \in \mathrm{supp}(\rho(s))
\text{ for some } s \in \mathcal{S},
\text{ and } s' \notin \pi^{-1}(\pi(\mathcal{S}
\setminus \mathcal{A}))\}.
\]
This is the set of states produced by the
repair that were invisible within the projection
used to design it.
\end{definition}

\section{Timeline of Model Construction and Collapse}

The following table maps Maxim's successive
interpretive frameworks against the events that
generate and destroy each one.

\medskip
\begin{center}
\begin{tabular}{p{3cm} p{4cm} p{4.5cm}}
\toprule
\textbf{Framework} & \textbf{Generating Event} &
\textbf{Collapsing Event} \\
\midrule
Robinson Crusoe / explorer & Crash landing on Saraksh &
Encounter with forced labor camp and police \\
\addlinespace
Totalitarian state & Observation of political institutions,
surveillance, compulsory ideology &
Discovery of broadcast mechanism and its scope \\
\addlinespace
Broadcast control model & Revelation from underground contact &
Encounters with mutants, Wanderer; consequences of tower attack \\
\addlinespace
Liberation project & Commitment to destroying Control Center &
Wanderer's account of cascade consequences \\
\addlinespace
Post-liberation stabilization & Novel's ending & Remains
unresolved; framework not yet constructed \\
\bottomrule
\end{tabular}
\end{center}
\medskip

Each row represents a projection that successfully
explains the evidence available at the time of
its adoption and fails when the next level of
hidden structure becomes visible. The table
illustrates that the collapses are not random
but follow the hierarchical architecture of
concealment described in Part~I.

% ---------------------------------------------------------------
% BIBLIOGRAPHY
% ---------------------------------------------------------------
\newpage
\begin{thebibliography}{99}

\bibitem{strugatsky1969}
Arkady Strugatsky and Boris Strugatsky.
\textit{Prisoners of Power} [\textit{Obytaemyy
Ostrov}]. Originally published in \textit{Neva},
1969. English translation by Helen Saltz Jacobson.
New York: Macmillan, 1977. Second English
translation by Andrew Bromfield. Chicago: Chicago
Review Press, 2014.

\bibitem{strugatsky1964}
Arkady Strugatsky and Boris Strugatsky.
\textit{Hard to Be a God} [\textit{Trudno byt'
bogom}]. Moscow: Molodaya Gvardiya, 1964.
English translation by Olena Bormashenko.
Chicago: Chicago Review Press, 2014.

\bibitem{strugatsky1979}
Arkady Strugatsky and Boris Strugatsky.
\textit{Beetle in the Anthill} [\textit{Zhuk
v muraveinike}]. Moscow: Molodaya Gvardiya, 1980.

\bibitem{postman1985}
Neil Postman. \textit{Amusing Ourselves to Death:
Public Discourse in the Age of Show Business}.
New York: Viking Penguin, 1985.

\bibitem{gatto2002}
John Taylor Gatto. \textit{Dumbing Us Down:
The Hidden Curriculum of Compulsory Schooling}.
Gabriola Island: New Society Publishers, 2002.

\bibitem{wiener1948}
Norbert Wiener. \textit{Cybernetics: Or Control
and Communication in the Animal and the Machine}.
Cambridge: MIT Press, 1948.

\bibitem{ashby1956}
W. Ross Ashby. \textit{An Introduction to
Cybernetics}. London: Chapman and Hall, 1956.

\bibitem{simon1969}
Herbert A. Simon. \textit{The Sciences of the
Artificial}. Cambridge: MIT Press, 1969.

\bibitem{popper1963}
Karl Popper. \textit{Conjectures and Refutations:
The Growth of Scientific Knowledge}. London:
Routledge and Kegan Paul, 1963.

\bibitem{kuhn1962}
Thomas S. Kuhn. \textit{The Structure of
Scientific Revolutions}. Chicago: University of
Chicago Press, 1962.

\bibitem{foucault1975}
Michel Foucault. \textit{Discipline and Punish:
The Birth of the Prison}. Paris: Gallimard, 1975.
English translation by Alan Sheridan. New York:
Pantheon Books, 1977.

\bibitem{scott1998}
James C. Scott. \textit{Seeing Like a State:
How Certain Schemes to Improve the Human
Condition Have Failed}. New Haven: Yale
University Press, 1998.

\bibitem{hayek1945}
Friedrich A. Hayek. ``The Use of Knowledge in
Society.'' \textit{American Economic Review}
35, no.~4 (1945): 519--530.

\bibitem{ostrom1990}
Elinor Ostrom. \textit{Governing the Commons:
The Evolution of Institutions for Collective
Action}. Cambridge: Cambridge University Press,
1990.

\bibitem{sen1999}
Amartya Sen. \textit{Development as Freedom}.
New York: Anchor Books, 1999.

\bibitem{holland1995}
John H. Holland. \textit{Hidden Order: How
Adaptation Builds Complexity}. Reading:
Addison-Wesley, 1995.

\bibitem{kauffman1993}
Stuart A. Kauffman. \textit{The Origins of
Order: Self-Organization and Selection in
Evolution}. Oxford: Oxford University Press,
1993.

\bibitem{pearl2000}
Judea Pearl. \textit{Causality: Models,
Reasoning, and Inference}. Cambridge: Cambridge
University Press, 2000.

\bibitem{lewis1973}
David Lewis. \textit{Counterfactuals}.
Cambridge: Harvard University Press, 1973.

\bibitem{stalnaker1968}
Robert Stalnaker. ``A Theory of Conditionals.''
In \textit{Studies in Logical Theory}, edited
by Nicholas Rescher, 98--112. Oxford: Blackwell,
1968.

\bibitem{flyxion2026rsvp}
Flyxion. \textit{Axioms for a Falling Universe:
A Unified Field Theory of the Relativistic
Scalar--Vector Plenum}. Independent Researcher,
May 2026. Available at:
\url{https://standardgalactic.github.io/alphabet/projects/axioms-for-a-falling-universe.pdf}

\bibitem{flyxion2026hydra}
Flyxion. \textit{HYDRA: Architecture, Geometry,
and the Problem of Coherent Agency}. Independent
Researcher, 2026. Available at:
\url{https://standardgalactic.github.io/playfloor/epistemology/hydra_architecture.pdf}

\bibitem{flyxion2026hiddencurvature}
Flyxion. \textit{Hidden Curvature: Constraint
Geometry, Projection, and the Reconstruction
of Meaning}. Independent Researcher, 2026.
Available at:
\url{https://standardgalactic.github.io/playfloor/admissibility/admissibility-field.pdf}

\bibitem{flyxion2026repair}
Flyxion. \textit{Repair as a Fundamental
Category: Toward an Ontology of Restoration}.
Independent Researcher, 2026. Available at:
\url{https://standardgalactic.github.io/alphabet/research/repair_as_fundamental.pdf}

\end{thebibliography}

\end{document}
